\chapter{函数}

微积分研究的对象主要是函数。其主要的动机并不是思辨的需要，而时常是更实际的经济利益的驱动。例如三角学，由古希腊就有系统的研究，这是为了研究天体的运行，也是为了航海的需要。所以当时不只是有平面三角学还有球面三角学。随着资本主义时代的来临，人类活动的广度与深度大为增加，例如对数的出现显然是为了简化航海中的计算。值得注意的是，这些活动甚至直接影响到学术界的最上层。下一章开始时我们会讲到牛顿，还有哈雷对天体运动的研究与航海的关系。现在我们将指出，对数的一位早期创始者尼古拉·麦卡脱\footnote{Nicholas Mercator，1620--1687。数学史上有两位麦卡脱，另一位是 Gerhard Kremer Mercator，1512--1594，是著名的地图制作者，现在地图上使用的麦卡脱投影法就是他创造的。他为制作地图所引进的一些名词，现在我们仍然使用在微分流形理论中，见第七章。这里一位的姓名中 Mercator 一字在德文文献中常作 Kaufman，商人之意。其实 Mercator 就是商人一词的拉丁文说法，英文的 Merchant 也来自此词，所以可以怀疑是否二人真的姓 Mercator，或者应译为“某某老板”？可以肯定的是，他们都不是职业数学家，而从事商贸。}实质上是用积分
\[
\int_0^x\frac{dx}{1+x}
\]
来定义对数。牛顿接受了他的思想，并把自己最得意的二项级数用上，与麦卡托互相独立地得到了 $\log(1+x)$ 的幂级数展开式。然后又用非常精巧的反演方法，得到了 $e^x$ 的幂级数展开式。牛顿也用同样的方法来研究
\[
\arcsin x=\int_0^x\frac{dx}{\sqrt{1-x^2}}
\]
以及它的反函数 $\sin x$，一直到 1715，布鲁克·泰勒（Brook Taylor，1685--1731）\jiaokan{原文作“Robert Taylor”。通常所称泰勒公式由英国数学家 Brook Taylor（1685--1731）系统发表，今据通行史料改正。}发表了“增量法及其逆”一文，才给出了我们现在知道的泰勒公式。所以，历史的发展与我们现在的教科书有些矛盾，先有反函数 $\log(1+x)=\int_0^x\frac{dx}{1+x}$ 和 $\arcsin x=\int_0^x\frac{dx}{\sqrt{1-x^2}}$ 的级数展开式，后来才有 $e^x$ 和 $\sin x$ 的级数展开式。为什么现在通用的微积分教本要用与历史不符的讲法呢？一个重要原因是教学的需要。首先看到这一点的欧拉，把关于微积分基础的讲授放在前面作为预备知识，于是多年以后三角学完全在中学里了。大学里先是给出了
\[
e=\lim_{n\to\infty}\left(1+\frac1n\right)^n,
\]
作为有界上升数列必有极限的一个重要的例子。可是为什么重要？为什么以 $e$ 为底的对数称为“自然”？都只好不讲了。$e^{i\theta}=\cos\theta+i\sin\theta$ 按逻辑的要求应该放到另一门课去讲，因为必须先定义什么是复变量的幂级数（其实这只是一句话的问题）。现在，我们既然假设读者都学过微积分，就有可能用一种比较接近历史真实的、比较接近物理学的讲法。我们的出发点是：一个函数 $y=y(x)$ 可以作为一个微分方程的柯西问题
\[
\left\{
\begin{aligned}
\frac{dy}{dx}&=f(x,y),\\
y(x_0)&=y_0
\end{aligned}
\right. \tag{A}
\]
之解来定义。这里 $y$ 可以是一个函数，也可以是一个向量 $y=(y_1,\ldots,y_n)$（当然 $f$ 也相应地成为一个向量 $f=(f_1,\ldots,f_n)$）。其实欧拉以及现在许多深一点的微积分教本都是这样作的，特别当涉及复变量时是如此。那么它是否能定义一个（而不是许多个）函数（也不是什么都没有定义）呢？从数学上讲这就是常微分方程的柯西问题的存在与唯一性定理。为此下面的条件就足够了：

\textbf{常微分方程的柯西问题（A）的存在与唯一性定理}\quad 设（1）右方的向量值函数 $f(x,y)$ 的每一个分量均定义在 $\Omega$ 中，$(x_0,y_0)\in\Omega$ 而且在 $\Omega$ 的子集
\[
|x-x_0|\le a,\qquad |y-y_0|\le b
\]
中对 $(x,y)$ 连续，对 $y$ 满足利普希茨条件，则（A）必有唯一解存在于 $[x_0-h,x_0+h]$ 中。这里
\[
h=\min\{a,b/M\},\qquad
M=\sup_{|x-x_0|\le a,\,|y-y_0|\le b}|f(x,y)|.
\]

这个定理我们就不来证明了。不过希望读者注意，这是一个很细致的定理，在实际用起来的时候一些似乎是很不起眼的地方，例如 $h$ 的大小，都会出问题。

我们上面已经说了欧拉就这样来研究过一些函数。但是欧拉的时代并没有这样的定理。这倒不是因为这个定理的证明很难，而是因为当时人们并不认为需要证明这样的定理，而且证明这个定理时必不可少的收敛性的概念，欧拉时代的数学家并不明白。例如看一个质点的运动，按牛顿第二定律，运动方程就是一个常微分方程组，如果给出了一定的初始条件，则该质点的运动状况就完全确定了。而且欧拉的想法是：把时间分成一小段一小段。在第一小段（从 $x=x_0$ 到 $x=x_0+h$），设 $f(x,y)=f(x_0,y_0)$，于是得到
\[
y=y_0+f(x_0,y_0)(x-x_0),
\]
这样到 $x=x_1=x_0+h$ 时，$y=y_1=y_0+f(x_0,y_0)h$。在第二小段时间，从 $x=x_1$ 到 $x=x_1+h$ 中不妨用 $f(x_1,y_1)$ 代替 $f(x,y)$，而得
\[
y=y_1+f(x_1,y_1)(x-x_1).
\]
仿此进行下去，我们会得到一条折线，称为欧拉折线。欧拉时代的人们认为 $h\to0$ 时就会得到真实的运动状况，而没有想到收敛性是要证明的。所以到 18 世纪为止，人们实际上是以物理的分析代替数学的分析。而且从哲学上看，这样的想法没有任何问题，因为世界上的一切都服从某种决定论，存在与唯一性定理无非是哲学上的决定论的数学表示。哲学上对了，数学上也就犯不着再自寻烦恼了。顺便说一下，不但欧拉折线是欧拉的创造，连牛顿的运动定律即以写成微分方程（A）也始自欧拉；所以我们常讲的牛顿运动方程也不妨称为牛顿--欧拉方程。下一章 §6 我们还要分析这里涉及的方法论问题，但是到了 19 世纪就发现这样做在数学上会出毛病。首先提出要证明存在与唯一性的是柯西。他在 1820--1830 年间发现仅有 $f(x,y)$ 的连续性并不足以保证唯一性，于是他补充上“$f_y(x,y)$ 也连续”的条件。利普希茨（Lipschitz，1832--1903）在 1876 年把这个条件减弱为我们现在熟知的利普希茨条件。人们还找到了不符合这类条件就会破坏唯一性的反例。而为了证明存在性，则 $f(x,y)$ 的连续性就够用了。

我们的读者都受到过相当程度的逻辑推理的训练，当然会问：这样做会不会是循环论证？肯定不会。在证明上述定理时我们没有用到任何的具体函数——不论是指数对数还是三角函数、反三角函数。还会问：这样做岂非杀鸡用牛刀？我们重新回到历史走过的老路，是为了看一下，微积分的发展怎样帮助我们更深刻地认识大自然的规律。从这个意义上讲，指数函数当然不是“鸡”，说它是“龙”也不为过。总之这不只是特例，而是十分重要的数学模型。

\section{增长的数学模型：指数与对数}

\textbf{1. 指数函数及其基本性质}\quad 在自然界和人类社会中有许多量的增长之快慢与该量现在的值成正比。例如人在银行的存款，在利率一定的条件下，存款增长多少，亦即利息之多少是与现在存款量（作为本金）成正比的。这里存款计息应该是按一定时期（年、半年）来计算的。在计息期末计息之前，存款额我们约定按期初之值计算，即是本金。又如某种细菌，某一个菌株是否恰好在一个时刻分裂可以是随机的，但在一定条件下，在一段时间中细菌群体的一个百分比（设为 $\lambda$）总是在分裂的，细菌就这样繁殖。某一定量的放射性物质的蜕变也是这样。对这种现象，既可以把它作为连续变化过程来处理（例如细菌的分裂，放射性元素原子的蜕变），也可以把它作为离散现象来处理。例如银行存款的增长是离散的，即使存期不到，也是按天计算，不会对时间连续计算。但是下面我们会看到，即使这个问题也可以作为连续变化的量来看待。总之，对这类问题可以从连续变量与离散变量两个角度来讨论。下面我们首先从连续变量的角度来讨论。

从连续变化的角度看，上述问题的数学模型是
\[
\left\{
\begin{aligned}
\frac{dy}{dt}&=\lambda y, &&\text{(1)}\\
y(0)&=y_0. &&\text{(2)}
\end{aligned}
\right.
\]
这个问题可以化简到 $\lambda=1$（因为只需作一个变换，即引入新的自变量 $t_1=\lambda t$ 即可）的情况，又可以假设 $y_0=1$。因为若得到 $y_0=1$ 的解 $E(t)$，则 $y_0E(t)$ 就是问题（1）、（2）之解。从上面讲到的存在与唯一性定理立即可知这时（1）、（2）有解 $E(t)$ 存在，而且 $E(t)$ 在其定义域上是连续可微函数。从我们熟知的常微分方程知识立即可知
\[
E(t)=e^t. \tag{3}
\]
这里 $e=\lim_{n\to\infty}(1+1/n)^n$，但因我们目前还没有定义 $e$，所以暂时还不知道（3）式是什么意思。但是我们有极为重要的

\textbf{定理 1（加法定理）}\quad 上述 $E(t)$ 之定义域是 $\mathbf R:-\infty<t<+\infty$，且在其定义域内有
\[
E(t_1+t_2)=E(t_1)E(t_2),\qquad t_1,t_2\in\mathbf R. \tag{4}
\]

\textbf{证}\quad 用读者已熟悉的初等方法自然可得以上一切结论，但是我们现在不能用（3）式，而要回到微分方程（1）。如果用初等方法求积（1）（令 $\lambda=1$）
\[
\frac{dy}{y}=dt,
\]
双方对 $t$ 求积分，就应考虑到如果 $y=y(t)$ 在某个 $t$ 值处为 $0$ 怎么办？所以严格的证法应该分别两种情况：一是 $y(t)$ 处处不为 $0$，这时可用求积法来处理它；但要证明相应于这个 $y(t)$ 的 $t$ 值充满 $\mathbf R$ 并非易事。二是 $y(t)$ 在某处为 $0$。所以，初等的求积方法其实不是一个好方法。我们还可以应用读者可能已学过的一个解的存在唯一性定理。在那里，我们证明了解在 $[x_0-h,x_0+h]$ 中存在，而 $h=\min\{a,b/M\}$。但是这里也有困难，因为 $a,b$ 固然可以取得任意大，$M$ 也会随之增大。所以在一般的常微分方程解的存在与唯一性定理后，专门有结果指出，对于线性方程（我们用矩阵形式来表示）
\[
\frac{dy}{dt}=A(t)y,
\]
若 $A(t)$ 在区间 $[a,b]$ 上连续，则解的定义域必包含 $[a,b]$，所以对于常系数线性方程，解的定义域必为 $\mathbf R$。这一切都属于比较细微的地方而可能会注意不到。我们下面要利用常系数线性方程的平移不变性来证明它；加法定理本身也就是平移不变性的表现。

设 $E(t)$ 已定义在 $[t_0-h,t_0+h]$ 中，$t_0=0$，于是令 $t_1=t_0+h_1$，$h_1=h\cdot\varepsilon$，$\varepsilon$ 是任意小正数，并记 $E(t_1)=E_1$，于是 $E(t)$ 在 $[t_0,t_0+h]$ 上也是柯西问题
\[
\frac{dy}{dt}=y,\qquad y(t_1)=E_1
\]
之解。令 $t=\tau+t_1$，则此问题又可化为
\[
\frac{dy}{d\tau}=y,\qquad y\big|_{\tau=0}=E_1.
\]
由前所述，其解为 $E_1E(\tau)$，而因为新的柯西问题形状与（1）、（2）相同，故由解的唯一性定理，在 $[t_1,t_1+h]$ 上，$E(t)=E_1E(\tau)=E(t_1)E(t-t_1)$，与（1）、（2）完全相同，所以这个解应可以定义于 $\tau\in[0,h]$ 上，即是说 $E(t)$ 不仅可以定义在 $[t_0,t_0+h]$ 上，而且可以定义在 $[t_1,t_1+h]=[t_0+h,t_0+h+h_1]$ 上，再由 $\varepsilon$ 之任意性，知 $E(t)$ 可以定义到 $[t_0+h,t_0+2h]$ 上。仿此进行下去知 $E(t)$ 可以定义到 $[t_0,+\infty)$ 上，同理也可以定义到 $(-\infty,t_0]$ 上。总之 $E(t)$ 之定义域为 $(-\infty,+\infty)$。

由前面已证的
\[
E(t)=E_1E(\tau)=E(t_1)E(t-t_1),
\]
令 $t=t_1+t_2$，由于 $E(t)$ 之定义域已知是整个 $\mathbf R$，所以 $t_1,t_2$ 取任意实数值均是允许的。这样，上式就给出了加法定理：对任意实数 $t_1,t_2$ 均有
\[
E(t_1+t_2)=E(t_1)E(t_2).
\]
定理证毕。

上面证明的关键在于引入了新的自变量 $\tau$ 以后，柯西问题的形状不变，因而原来的定理现在仍然适用，这就是平移不变性。如果（1）表示某种规律，则当 $t$ 表示时间时，平移不变性就表示这个规律在任何时间都相同；如果 $t$ 表示空间位置，平移不变性就表示这个规律在各地都一样，常系数的微分方程就表示这类规律，因此是最重要的。

下面要问，$E(t)$ 究竟是什么函数？准确一些说，应该问，$E(t)$ 与我们过去熟知的那些函数有何关系？下面要证明 $E(t)$ 是指数函数。

1. 首先注意到 $E(t)$ 恒不为 $0$，因为若 $E(t_0)=0$，不妨设 $t_0=0$（令 $t=\tau+t_0$，再用平移不变性），则有 $dE(\tau)/d\tau=E(\tau),E(0)=0$，由解的唯一性定理应有 $E\equiv0$，但因 $E(0)=1$，所以 $E(t_0)=0$ 不能成立，从而 $E(t)$ 恒不为 $0$。$E(t)$ 既是连续函数，故知 $E(t)>0$。

2. $E(t)$ 既然定义在 $\mathbf R$ 上，则它把 $\mathbf R$ 映到什么地方？因为 $dE/dt=E>0$，所以 $E$ 是上升函数，它必把 $\mathbf R$ 单调地映到某区间 $(E(-\infty),E(+\infty))$ 上，所以要问 $E(\pm\infty)=?$ 由于
\[
\int_1^E\frac{dE}{E}=\int_0^t dt,
\]
令右方的 $t\to+\infty$，则得 $\int_1^E dE/E$ 在 $t\to+\infty$ 时趋于 $+\infty$，为此需要 $E\to+\infty$，所以 $E(+\infty)=+\infty$。当 $t\to-\infty$ 时，因为 $\int_1^E dE/E$ 在 $E=0$ 处发散，易见这时 $E\to0$，即是说 $E(-\infty)=0$。所以 $E(t)$ 把 $\mathbf R$ 单调地映为 $(0,+\infty)$。下面讨论 $E(t)$ 的反函数时，这个性质是其基础。

3. 最后证明 $E(t)=[E(1)]^t$。首先设 $t$ 是正整数，例如 $t=2$，
\[
E(2)=E(1+1)=E(1)\cdot E(1)=E(1)^2.
\]
同理 $E(n)=[E(1)]^n$，$n$ 为正整数。其次看 $E(-n)$，$n$ 仍为正整数。注意到由定义 $E(0)=1$（就是（2）式），所以
\[
1=E(0)=E(n+(-n))=E(n)\cdot E(-n)=[E(1)]^nE(-n),
\]
而 $E(-n)=[E(1)]^{-n}$。（注意前面已证明 $E(t)$ 恒不为 $0$，故 $E(1)\ne0$，这样，上面的推导才有意义。）总之
\[
E(n)=[E(1)]^n,\qquad n=0,\pm1,\pm2,\ldots.
\]
再看 $t=n/m$ 为有理数（$m,n$ 为整数，设 $m>0$）的情况。首先，由 $1=1/m+\cdots+1/m$（$m$ 项），有
\[
E(1)=E\left(\frac1m+\cdots+\frac1m\right)=\left[E\left(\frac1m\right)\right]^m,
\]
\[
E\left(\frac1m\right)=[E(1)]^{1/m}.
\]
于是对 $n>0$，$n/m=1/m+\cdots+1/m$（$n$ 项），有
\[
E\left(\frac nm\right)=\left[E\left(\frac1m\right)\right]^n=[E(1)]^{n/m}.
\]
而 $n<0$ 时，则利用 $E\left(\frac nm+\frac{-n}{m}\right)=E(0)=1$，从而
\[
E\left(\frac nm\right)=\left[E\left(\frac{-n}{m}\right)\right]^{-1}=[E(1)]^{n/m}.
\]
（注意 $-n>0$。）所以 $t$ 为有理数时 $E(t)=[E(1)]^t$ 也成立。

最重要的一步是 $t$ 为一般实数的情况。在通常的微积分教本中定义 $a^t$ 时比较马虎一些的就说 $a^t$ 是连续函数。（对这句话则不加论证，而且 $a^t$ 尚未定义又怎么知道 $a^t$ 是连续函数呢？）在比较严肃一些的教本中都是取一串有理数 $t_n\to t$，然后讨论 $\lim a^{t_n}$，当 $a\ne1$（现在不会出这样的事，因为 $E(t)$ 是上升函数而于 $t>0$ 时，$E(t)>E(0)=1$）例如可以取一串上升的 $t_n\to t$，再用单调有界数列定理证明这个极限存在。同时还要证明，不论 $t_n$ 如何取，不论它是否单调，$\{a^{t_n}\}$ 都有相同的极限，用这个极限值作为 $a^t$ 之定义，这样才把指数函数 $a^t$ 定义在整个 $\mathbf R$ 上。总之是用上了实数理论的全武行！本书第六章会详细讲实数理论，而且为了以后考虑一般函数空间，是用柯西序列来定义实数的。但就现在的情况，则可以完全跳过去。因为，由存在与唯一性定理已经知道 $E(t)$ 在整个实数轴 $\mathbf R$ 上有定义而且连续，所以哪怕 $t$ 是一般实数，$E(t)$ 已有定义而且在 $t$ 处连续，根本用不着证明。何况 $E(t_n)=[E(1)]^{t_n}$ 已经验证，所以对一般的实数 $t$ 即有
\[
E(t)=\lim_{n\to\infty}E(t_n)=\lim_{n\to\infty}[E(1)]^{t_n}.
\]
因为 $E(t)$ 在 $\mathbf R$ 的一个稠密集（有理数集）上与初等方法定义的 $[E(1)]^t$ 相同，因此在整个 $\mathbf R$ 上都不妨用 $[E(1)]^t$ 这样的记号来写，并且称之为指数函数。总之我们看到，我们回避了对无理数指数定义 $a^t$ 并研究其性质这个困难问题，而变成了对一个由柯西问题（1）、（2）所定义的函数证明其在 $\mathbf R$ 的一个稠密子集（有理数集）上与初等方法定义的 $a^t$（$a=E(1)$）相一致的问题。至于 $E(t)$ 的性质如上面的 1、2 两点全是利用微分方程方法证明的。

这样一来岂非实数理论等等都是无用而多余的吗？不是。数学中任何一个重要问题总有其关键之处，用不同的解决方法关键点可能不同，但是它们困难的程度都是一样的，甚至这些解决方法本质上都是一样的。上面我们把问题归结为常微分方程解的存在与唯一性定理，而这个定理的证明又少不了证明某个函数序列的一致收敛性。抽象地说来这也是一个完备性问题，与实数理论要解决的问题是一样的，详见第六章度量空间一节，这种情况以后还会出现。

现在余下的问题是，我们很相信
\[
E(1)=\lim_{n\to\infty}\left(1+\frac1n\right)^n,
\]
可是怎么证明呢？我们的方法是回到反函数上去，并且用离散的方法处理它。

\textbf{2. 增长的离散模型}\quad 现在设自变量 $t$ 不是连续变化的，而是“跳动”着在变，其“步长”是一个有限数 $h>0$，即每一次都是 $t\mapsto t+h$，或写作 $\Delta t=h$，于是 $y(t)$ 也相应有改变量 $\Delta y=y(t+h)-y(t)$。现在设
\[
\frac{\Delta y}{\Delta t}=\lambda y,\qquad y(0)=y_0. \tag{5}
\]
问 $y(t)$ 是什么函数？（5）称为一个差分方程，它是上述增长问题的“离散模型”。人们时常会以为只有连续模型（1）、（2）才是“真实”的，而离散模型则只是近似的。这是不对的。例如 $y(t)$ 是某人在银行的存款，则后一个条件表明他在期初（$t=0$ 时）存入了 $y_0$，利率为 $\lambda$，则 $y(t)$ 就是他在 $t$ 时的本利和。这个问题本质上就是离散的。由（5），并以一个存期为单位，即令 $h=1$，有
\[
y(1)-y(0)=\lambda y(0),
\]
亦即
\[
y(1)=(1+\lambda)y(0).
\]
所以由 $y(0)$ 开始会得到
\[
y(n)=(1+\lambda)y(n-1)=\cdots=(1+\lambda)^ny(0).
\]
如果总的存期仍为 1，但把它分成 $n$ 个相等的小间隔，以 $1/n$ 为短的存期。因为存期变短，所以利率也应相应变小，设为 $\lambda/n$，仍然以存期的数目的 $1/n$ 为单位，则原来的一个存期现在要算作 $n$ 个小存期。仍用上面的结果可知，在长存期 1 之末，本利和应该用上式的 $y(n)$ 来计算而有
\[
y(n)=\left(1+\frac{\lambda}{n}\right)^ny(0). \tag{6}
\]
不过这时（6）给出的是一个长存期末的本利和。两个长存期末的本利和是
\[
\left(1+\frac{\lambda}{n}\right)^{2n}y(0),\ldots,
\]
$t$ 个长存期末的本利和是
\[
\left(1+\frac{\lambda}{n}\right)^{nt}y(0).
\]
令 $n\to\infty$，会得到连续的模型。用 $1/n$ 作为时间单位，即有
\[
y(t)=\lim_{n\to\infty}\left(1+\frac{\lambda}{n}\right)^{nt}y(0)=e^{\lambda t}y(0). \tag{7}
\]
把它和连续模型得到 $E(t)=[E(1)]^t$ 比较，更令人想到 $E(1)$ 应该就是
\[
e=\lim_{n\to\infty}\left(1+\frac1n\right)^n,
\]
但是严格地证明这一点是不容易的。我们不妨来看一看欧拉是怎样来研究这些问题的。首先，我们要给出一个“实际”计算 $E(t)$ 的方法，为此令
\[
E(t)=\sum_{n=0}^{\infty}a_nt^n,
\]
则由 $E(0)=1$ 应该有 $a_0=1$，又由 $E'(t)=E(t)$ 应该有
\[
\sum_{n=0}^{\infty}na_nt^{n-1}=\sum_{n=0}^{\infty}a_nt^n.
\]
比较同类项的系数有
\[
(n+1)a_{n+1}=a_n,\qquad a_{n+1}=\frac{1}{n+1}a_n,
\]
所以
\[
a_{n+1}=\frac{1}{n+1}a_n=\frac{1}{(n+1)n}a_{n-1}=\cdots=\frac{1}{(n+1)!}a_0=\frac{1}{(n+1)!},
\]
所以
\[
E(t)=1+\frac{t}{1!}+\frac{t^2}{2!}+\cdots+\frac{t^n}{n!}+\cdots. \tag{8}
\]
这是一个形式解。当然考虑其收敛性并不困难，所以我们认为已证明了（8）在 $-\infty<t<+\infty$ 上的收敛性以及许可在其上逐项微分、逐项积分。其实对复数 $t$ 也是一样的，详见下一章 §4。但是在欧拉的时代，人们对收敛性问题的重要性还很不注意，真正严重地注意到收敛性的数学家应该提到高斯和阿贝尔，即已经是 19 世纪初年了。欧拉的伟大在于他对数学公式的推演有非凡的才能，对正确的结论有出乎常人的洞察力，虽然自己不甚严格，但得出了许多极为重要的结果。即以 $e$ 的问题为例，他先用二项定理对正整数 $n$ 得出
\[
\left(1+\frac1n\right)^n=1+\frac n1\frac1n+\frac{n(n-1)}{2!}\frac1{n^2}+\cdots+\frac{n!}{n!}\frac1{n^n},
\]
令 $n\to\infty$，欧拉对每一个有限项取极限，则右方的前若干项成为
\[
1+\frac1{1!}+\frac1{2!}+\cdots+\frac1{n!}.
\]
再令项数趋于无穷，有
\[
e=1+\frac1{1!}+\frac1{2!}+\cdots+\frac1{n!}+\cdots. \tag{9}
\]
欧拉的作法从今天看来当然是不能接受的。但是他不但得到了正确的（9）式，而且他用类似的方法处理 $\left(1+\frac1n\right)^{nt}$，得到了
\[
e^t=1+\frac t{1!}+\frac{t^2}{2!}+\cdots+\frac{t^n}{n!}+\cdots, \tag{10}
\]
这与（8）完全一样。所以我们立即有 $E(1)=e$。可惜，这个证明是不严格的。进一步我们回到 $E(t)$ 的反函数。

\textbf{3. 对数函数}\quad $y=E(t)$ 有反函数吗？上面我们已经证明了 $E(t)$ 是单调函数，它把 $\mathbf R$ 单调地（从而是一对一的）、连续可微地映到正实数轴 $\mathbf R_+=\{x:x>0\}$ 上，所以，它必有一个反函数，即以 $E(1)$ 为底的对数 $\log_{E(1)}y$。我们把它写成 $t=L(y)$，定义在 $y>0$ 上而且把 $y>0$ 单调地、连续可微地映为 $-\infty<t<+\infty$，而且因为 $t=0$ 时 $y=E(0)=1$，所以 $L(1)=0$。由反函数的导数的定义知道
\[
\frac{dt}{dy}=\frac1y,\qquad t\big|_{y=1}=0.
\]
所以
\[
t=L(y)=\int_1^y\frac{dy}{y}. \tag{11}
\]
如果令 $y=1+x$，由上式经变量变换有
\[
t=L(1+x)=\int_0^x\frac{dx}{1+x}. \tag{12}
\]
但是，通常的微积分教本告诉我们，（11）恰好就是以 $e$ 为底的对数 $\ln y$，所以
\[
L(y)=\ln y. \tag{13}
\]
当 $y=E(1)$ 时，由反函数的定义知（11）中的 $t=1$，代入（13）有
\[
1=\ln E(1).
\]
从而 $E(1)=e$。但是这样做还不能令我们满意，因为我们都熟知
\[
e=\lim_{n\to\infty}\left(1+\frac1n\right)^n,
\]
怎样把 $E(1)$ 与这个极限式联系起来呢？我们已经看到，若用（1）、（2）的离散模型，会见到 $\left(1+\frac1n\right)^n$。离散模型的“极限”就是连续模型。这句话一般地说相当有“说服力”，但是总还欠严格。那么，能否证明 $E(1)$ 直接就是 $\lim\left(1+\frac1n\right)^n$ 呢？这是可以的，我们将在下面讨论。

现在我们要利用积分公式（11）直接讨论 $L(y)$ 的性质。首先是加法定理。对于 $y_1,y_2\ne0$，有
\[
L(y_1y_2)=\int_1^{y_1y_2}\frac{dy}{y}=\int_1^{y_1}\frac{dy}{y}+\int_{y_1}^{y_1y_2}\frac{dy}{y}.
\]
在第二个积分中作变换 $y=y_1\xi$，因为 $y_1\ne0$，这个变换是合法的，所以有
\[
L(y_1y_2)=\int_1^{y_1}\frac{dy}{y}+\int_1^{y_2}\frac{d\xi}{\xi}=L(y_1)+L(y_2).
\]
这就是对数的加法定理。本来它由 $L(y)=\log_{E(1)}y$ 以及对数的定义可以直接证明，但是用积分证明似更直接。

其次是两个基本的极限。在通常的微积分教本中，证明 $\frac{d}{dx}e^x=e^x$ 以及 $\frac{d}{dy}\ln y=1/y$ 要用两个基本的极限
\[
\lim_{x\to0}\frac{e^x-1}{x}=1,
\qquad
\lim_{x\to0}\frac{\ln(1+x)}{x}=1.
\]
其证明相当不易，但是用我们现在的方程，这两个不等式都成了几乎不证自明的事。

首先我们知道 $E'(t)=E(t)$，而且 $E(0)=1$，所以 $E'(0)=E(0)=1$，由导数的定义立即有
\[
E'(0)=\lim_{x\to0}\frac{E(x)-E(0)}x=\lim_{x\to0}\frac{e^x-1}{x}=1,
\]
因为我们已经知道了 $E(1)=e$（尽管还得严格证明），故 $E(x)=[E(1)]^x=e^x$。关于第二个极限式，利用 $L(y)=\log_{E(1)}y=\ln y$ 以及（12）式用积分中值定理
\[
\ln(1+x)=\int_0^x\frac{dx}{1+x}=\frac{x}{1+\xi},\qquad 0\le\xi\le x.
\]
令 $x\to0$ 即知 $\xi\to0$，而得
\[
\lim_{x\to0}\frac{\ln(1+x)}x=\lim_{\xi\to0}\frac1{1+\xi}=1.
\]

最后是幂级数展开式。对 $E(x)=e^x$，我们已得到处处收敛的（8）式；对于对数函数，当 $|x|<1$ 时，由于
\[
\frac1{1+x}=1-x+x^2-x^3+\cdots,
\]
逐项积分后，立即有，当 $|x|<1$ 时，
\[
\ln(1+x)=\int_0^x\frac{dx}{1+x}=x-\frac{x^2}{2}+\frac{x^3}{3}+\cdots+(-1)^{n-1}\frac{x^n}{n}+\cdots. \tag{14}
\]
这就是尼古拉·麦卡托当年的作法。

\textbf{4. 对数的离散处理——$E(1)=\lim_{n\to\infty}(1+1/n)^n$ 的证明}\quad 从历史上看，对数的出现要早于指数。上文说的欧拉的贡献都见于他的《无穷小分析引论》一书，出版于 1748 年，但纳皮尔（John Napier，1550--1617）在 1614 年就造出了第一本对数表，远早于牛顿的微积分。另一个瑞士人比尔吉（Jobst Bürgi，1552--1632）也在 1620 年发表了另一个对数表。对数的出现确实直接来自航海的需要，因为航海要作许多乘除法计算，当然人们希望用更简单的加减法来代替它们。这样就想到了用对数 $y=\log_b t$，它是指数函数 $t=b^y$ 的反函数。对数的底 $b$ 从一开始并不是 10，更不是 $e$，而是看如何选取计算最为简单而定。于是问题就成了给定一定的 $t$ 如何求相应的 $y$，这当然是很困难的。于是纳皮尔和比尔吉不约而同地想到：把问题反过来做，即给出许许多多的 $y$，再用幂运算来计算 $b^y$，得到许许多多的 $t$。他们希望这些 $t$ 尽量密集，这样，当你有了一个 $t$ 以后，很有可能这个 $t$ 正好纳皮尔从某一个 $y$ 算出来的，至少很接近纳皮尔算过的值，于是你就得到了相应于这个 $t$ 的 $y$，或者说，就得到了这个 $t$ 的对数，至少是其近似值。但是用 $b^y$ 算 $t$ 也不简单，除非 $y$ 是整数。这样，既要 $y$ 是整数，因而 $\Delta y=1$，从而 $y$ 不能太密集，又要求 $t$ 很密集，这个矛盾如何解决呢？办法是取 $b$ 尽可能接近 $1$（但 $b\ne1$，否则不论 $y$ 是什么值，$t=1$ 总是 $1$）。所以比尔吉取 $b=1+10^{-4}$，纳皮尔取 $b=1-10^{-7}$，纳皮尔用小于 $1$ 的 $b$ 作底是因为航海中需要计算的多是某角的正弦，因而小于 $1$。如用大于 $1$ 的 $b$ 作底，对数将是负数，用起来很不方便。

下面只限于讲比尔吉的作法。我们问当 $y$ 有一个增量 $\Delta y=1$ 时，$\Delta t$ 是多少？$t$ 的相对误差是
\[
\frac{\Delta t}{t}=\frac{(1+10^{-4})^{y+1}-(1+10^{-4})^y}{(1+10^{-4})^y}=\frac1{10^4}=\frac{\Delta y}{10^4},
\]
所以 $\Delta t$ 相当小，而 $t$ 的分布很密。上式其实就是差分方程
\[
\frac{\Delta y}{\Delta t}=\frac{10^4}{t}. \tag{15}
\]
如果用 $y/10^4=z$ 代替 $y$，则对数的定义式与上述差分方程成为
\[
\frac{\Delta z}{\Delta t}=\frac1t,
\]
\[
t=\left(1+\frac1{10^4}\right)^{10^4z}=\left[\left(1+\frac1n\right)^n\right]^z, \tag{16}
\]
这里 $n=10^4$，如果令 $n\to\infty$，则极限情况是
\[
\frac{dz}{dt}=\frac1t,
\]
即
\[
t=e^z.
\]
这就是 $e=\lim(1+1/n)^n$ 的真正的来源。可惜我们不知道欧拉是不是这样想的。但是欧拉对他发现的这个极限似乎情有独钟，所以把自己的姓氏用小写字母写上去了。当然，这不是“证明”，只是一种启发，要证明差分方程的解之极限确是微分方程之解，并非易事。证明见后文。

现在我们来求解差分方程，$y=0$（即 $z=0$）对应于 $t=1$，然后用相同的步长 $\Delta y=1$，即 $\Delta z=1/10^4$，从 $z_0=0$ 依次得到分点 $z_0<z_1<\cdots<z_N$，而相应的 $t$ 的分点是 $t_0=1<t_1<\cdots<t_N$。$\Delta t=t_i-t_{i-1}$ 是不相同的，但这并没有关系，因为由差分方程 $\Delta y/\Delta t=10^4/t$ 给出 $\Delta z/\Delta t=1/t$，亦即
\[
z_{i+1}-z_i=\frac1{t_i}\Delta t_i,
\qquad i=0,1,\ldots,N-1.
\]
把它们加起来，并且注意到 $z_0=0$，即有差分方程的解
\[
z_N=\sum_{i=0}^{N-1}\frac{\Delta t_i}{t_i}. \tag{17}
\]
显然它的极限即
\[
z=\int_1^t\frac{dt}{t}.
\]
因此，当 $\Delta z\to0$ 时不但差分方程变成微分方程 $dL(t)/dt=1/t$，而且差分方程的解（17）之极限也正是定义对数函数的积分。

现在我们可以来证明
\[
E(1)=\lim_{n\to\infty}\left(1+\frac1n\right)^n
\]
了。在通常的微积分教本中都要先证明此极限的存在性，并且记之为 $e$。我们的工作则不是证明 $e$ 的存在（证明见通常的教本），而是承认它的存在，只证明 $e=E(1)$。

证明的基本思想是从几何着眼。$y=E(t)$ 是微分方程 $dy/dt=y$ 的解，另外有差分方程（5）
\[
\frac{\Delta y}{\Delta t}=y. \tag{18}
\]
可是差分是那里算起呢？如图 2-1-1，在 $y=E(t)$ 上取两点 $P:(a,E(a))$，$Q:(a+h,E(a+h))$。为了在 $P,Q$ 之间考虑逼近 $y=E(t)$，如果要从 $P$ 开始，则用 $F(t)$ 逼近 $E(t)$，这里 $y=F(t)$ 的差分取为
\[
\frac{\Delta y}{\Delta t}=\frac{y-F(a)}{t-a}. \tag{18$_1$}
\]
也可以从 $Q$ 开始，用 $G(t)$ 逼近 $E(t)$，而 $y=G(t)$ 的差分取为
\[
\frac{\Delta y}{\Delta t}=\frac{y-G(a+h)}{t-a-h}. \tag{18$_2$}
\]
于是在（18$_1$）的情况，差分方程 $\Delta y/\Delta t=y$ 成为
\[
\frac{\Delta y}{\Delta t}=\frac{y-F(a)}{t-a}=F(a),
\quad\text{或}\quad
 y=F(a)+F(a)(t-a). \tag{19$_1$}
\]
在（18$_2$）的情况下则有
\[
\frac{\Delta y}{\Delta t}=\frac{y-G(a+h)}{t-a-h}=G(a+h),
\quad\text{或}\quad
 y=G(a+h)+G(a+h)(t-a-h). \tag{19$_2$}
\]

\begin{figure}[htbp]
\centering
\begin{tikzpicture}[bella figure,x=1.05cm,y=0.9cm,>=stealth,line cap=round,line join=round]
  \draw[->] (0,0)--(5.15,0) node[right] {$t$};
  \draw[->] (0,0)--(0,4.35) node[above] {$y$};
  \coordinate (P) at (1.45,1.45);
  \coordinate (Q) at (4.15,3.65);
  \coordinate (T) at (2.95,1.45);
  \draw[dashed] (P)--(1.45,0) node[below] {$a$};
  \draw[dashed] (Q)--(4.15,0) node[below] {$a+h$};
  % 按原图保持曲线、两条近似直线与交点的相对位置。
  \draw[thick] (P) .. controls (2.20,1.42) and (3.18,1.88) .. (Q);
  \draw[thick] (1.45,1.45)--(4.45,1.45);
  \draw[thick] (2.42,0.55)--(Q);
  \node[above left] at (P) {$P(a,E(a))$};
  \node[above right] at (Q) {$Q(a+h,E(a+h))$};
  \node[above] at (2.73,2.18) {$y=E(t)$};
  \node[right] at (4.28,1.45) {$y=F(t)$};
  \node[below left] at (2.56,0.83) {$y=G(t)$};
  \node[below right] at (T) {$T$};
\end{tikzpicture}
\caption*{图2-1-1}
\end{figure}

注意到 $E'(t)=E(t)$，所以 $E''(t)=E(t)$，但是前面我们已经说过 $E(t)>0$，所以 $y=E(t)$ 是一条向上凹的曲线，因此它必位于 $t=a$ 与 $t=a+h$ 两点之切线 $PT$ 与 $QT$ 之上侧（图 2-1-1）。如果我们记 $PT$ 为 $y=F(t)$，$F(t)$ 之表达式即（19$_1$），记 $QT$ 为 $y=G(t)$，$G(t)$ 的表达式为（19$_2$），所以我们有
\[
E(t)\ge F(t),\qquad E(t)\ge G(t), \tag{20}
\]
而且 $E(t)=F(t)$ 只在 $t=a$ 时成立，$E(t)=G(t)$ 只在 $t=a+h$ 时成立。

现在我们把 $t$ 轴上的区间 $[0,1]$ 作 $n$ 等分：$t_0=0<t_1<\cdots<t_n=1$，$t_i=i/n$。注意到 $E(0)=F(0)=1$，由（19$_1$），给出
\[
F(a+h)=F(a)+F(a)h=(1+h)F(a).
\]
于是令 $a=(i-1)/n,h=1/n$，从 $i=1$ 一直做到 $i=n-1$ 有
\[
\begin{aligned}
F(1)&=F\left(\frac nn\right)
=\left(1+\frac1n\right)F\left(\frac{n-1}{n}\right)
=\left(1+\frac1n\right)^2F\left(\frac{n-2}{n}\right)=\cdots\\
&=\left(1+\frac1n\right)^nF(0)
=\left(1+\frac1n\right)^nE(0)
=\left(1+\frac1n\right)^n.
\end{aligned}
\]
所以 $F(1)=\left(1+\frac1n\right)^n$ 实际上就是差分方程（18$_1$）之解。由于 $E(t)$ 的曲线在切线 $PT$ 上方，故
\[
E(1)>F(1)=\left(1+\frac1n\right)^n.
\]
用类似的方法来讨论 $G(t)$。上面我们求解（19$_1$）时，是用 $F(t)$ 来表示 $F(a+h)$，对于（19$_2$），则反过来要利用 $G(a+h)$ 来表示 $G(a)$，即由（19$_2$）令 $t=0$ 得出
\[
G(a)=G(a+h)-G(a+h)h=(1-h)G(a+h).
\]
现在我们从 $a=(n-1)/n$ 开始，而 $h=1/n$，注意到现在不是 $E(0)=F(0)=1$，而是 $G(1)=E(1)$，而不知道其值究竟是多少。用上式从 $a=(n-1)/n,h=1/n$ 开始逐次往前推算，有
\[
E(1)=G(1)=\left(1-\frac1n\right)^{-1}G\left(\frac{n-1}{n}\right)=\cdots=\left(1-\frac1n\right)^{-n}G(0).
\]
$E(t)$ 的图像也在切线 $G(t)$ 之图像上方。现在对 $t=0$ 应用这个事实（上面是在 $t=1$ 处应用 $E(t)>F(t)$ 的），有
\[
E(1)=\left(1-\frac1n\right)^{-n}G(0)<\left(1-\frac1n\right)^{-n}E(0)=\left(1-\frac1n\right)^{-n}.
\]
综合这两步，得到最后的不等式：
\[
\left(1+\frac1n\right)^n<E(1)<\left(1-\frac1n\right)^{-n}. \tag{21}
\]
这就与通常微积分教本中同时考虑两个序列 $\{(1+1/n)^n\}$ 和 $\{(1-1/n)^{-n}\}$ 完全一致了。不过在通常教材上先证明前一序列上升而有上界，后一序列下降而有下界，因而各有极限，并且证明这两个极限相等，记之为 $e$。现在我们则是承认这个事实，不过对它作了几何解释，即 $y=E(t)$ 的图像在任一区间 $[a,a+h]$ 上均在其左、右两端的两个切线上方，然后把 $[0,1]$ 分成 $n$ 个等分，并且对每个小区间 $[(i-1)/n,i/n]$ 应用上面的事实：或者对左端的 $PT$，我们从左向右，依次令 $i=1,2,\ldots,n$。在左端 $i=1$ 处用 $E(0)=F(0)=1$，而在右端的 $i=n$ 处应用 $E(1)>F(1)$ 得出 $E(1)>\left(1+1/n\right)^n$；或者对右端的 $QT$ 则由右向左，依次令 $i=n-1,\ldots,0$，并在 $i=n-1$ 处应用 $E(1)=G(1)$，在最左端的 $i=0$ 处应用 $1=E(0)>G(0)$ 而得出 $E(1)<\left(1-1/n\right)^{-n}$，总之有（21）式。

注意到
\[
\lim_{n\to\infty}\left(1+\frac1n\right)^n=e,
\]
\[
\begin{aligned}
\lim_{n\to\infty}\left(1-\frac1n\right)^{-n}
&=\lim_{n\to\infty}\left(1+\frac1{n-1}\right)^n\\
&=\lim_{n\to\infty}\left(1+\frac1{n-1}\right)^{n-1}
  \lim_{n\to\infty}\left(1+\frac1{n-1}\right)=e.
\end{aligned}
\]
立即有
\[
E(1)=e. \tag{22}
\]
这就是我们最后的目标。

\section{周期运动和三角函数}

\textbf{1. 匀速圆周运动}\quad 周期运动在自然界中几乎是无处不见。天体的运行，不论按日心说或者按地心说，其最简单的“组成成分”都是匀速的圆周运动。正因为日月五星都围绕着一个中心、周而复始，所以研究它们所必需的三角函数，在古希腊数学中占据了很重要的地位。与上一节讨论指数函数与对数函数一样，我们要从微积分学的角度来看待它。但是微积分学既然开始于力学
